\paragraph{Introduction.}
Isoperimetric inequalities, which relate the area to the length of its boundary, have been widely studied in mathematics.
For example, the isoperimetric inequality states that for a closed curve with length $L$, the area $A$ of the region satisfies $4\pi A \leq L^2$ in two-dimensional Euclidean space.

We study influence and isoperimetric inequalities on the Boolean hypercube $\Omega = \{\pm 1\}^n$ with the distribution $\mu$ over $\Omega$.
For a subset $S\subseteq \Omega$, the analogue of its volume is measured by $\mu(S)$.
In the following, we represent a subset $S\subseteq\Omega$ by a Boolean function $f:\Omega\mapsto \{\pm 1\}$ ($f^{-1}(+1) = S$).
Boolean functions naturally model voting rules, graph properties, and
set families.
For example, in a voting rule, each coordinate records an individual
vote, $f(x)$ gives the collective decision, and $S$ represents the
winning coalitions.
In this notation, $\Var[\mu]{f} = 4\mu(S)(1-\mu(S)) \geq 2\min\{\mu(S),1-\mu(S)\}$ serves as the analogue of a volume term, measuring the smaller side of the partition with fixed boundary.

There are several natural notions of boundary on the Boolean hypercube. Two classical ones are based on the \emph{influence} and the \emph{sensitivity} of coordinates.
We introduce the influence of the $i$-th coordinate and the sensitivity with respect to the Boolean function $f$ as follows:
$$
    \mathsf{Inf}_{\mu,i}(f) = \Pr[\mu]{f(x)\neq f(x^i)}
    \quad \text{and} \quad s_\mu(f)(x) = \sum_{i=1}^n \ind[f(x)\neq f(x^i)],
$$
where $x^i$ is obtained by flipping the sign of the $i$-th coordinate of $x$.
The influence $\mathsf{Inf}_{\mu,i}(f)$ measures the volume affected by flipping the $i$-th coordinate while the sensitivity counts how many coordinates change the membership for a given $x$. 
In our voting model, they correspond to the importance of the $i$-th voter and the number of voters who can individually change the outcome of an election, respectively.
The \emph{surface area}, $\mathsf{SA}_\mu(f):=\E[\mu]{\sqrt{s_\mu(f)}}$, is the boundary area based on square-root sensitivity.
For the uniform measure on the Boolean hypercube, $\mathsf{SA}_\mu(f)$ corresponds to the Houdr\'e--Tetali isoperimetric constant $\varphi_{1/2}$, which interpolates between edge conductance and vertex expansion \cite{lau2024houdr}.
These quantities have a natural interpretation on the hypercube graph.
For a subset of vertices $S=f^{-1}(+1)$, both the \emph{total influence} $\sum_{i=1}^n\mathsf{Inf}_{\mu,i}(f)$ and the surface area $\mathsf{SA}_\mu(f)$ quantify the size of the edge boundary $\partial S=\{(i,x):f(x)\neq f(x^i)\}$ from different perspectives.

Let us consider $\mu$ to be the uniform distribution.
One fundamental isoperimetric inequality is the Poincar\'e inequality: for Boolean function $f:\Omega\mapsto \{\pm 1\}$, $\Var[\mu]{f} \leq \sum_{i=1}^n \mathsf{Inf}_{\mu,i}(f)$, which guarantees that $\max_{i\in[n]} \mathsf{Inf}_{\mu,i}(f) \ge \Var[\mu]{f}/n$.
In this proposal, we focus primarily on the following inequalities, for any Boolean function $f:\Omega\mapsto \{\pm 1\}$:
\begin{itemize}
    \item [1.] \textbf{Kahn--Kalai--Linial (KKL) inequality \cite{KKL88}}: $\displaystyle \max_{i\in[n]} \mathsf{Inf}_{\mu,i}(f)\gtrsim \frac{\log n}n\Var[\mu]{f}$.
    \item [2.] \textbf{Talagrand inequality \cite{Talagrand93}}: $\displaystyle \mathsf{SA}_\mu(f)\gtrsim\Var[\mu]{f}\sqrt{\log \frac e{\Var[\mu]{f}}}$.
\end{itemize}
For a voting rule, KKL guarantees a voter with influence $\Omega((\log n)/n)$, a logarithmic improvement over the Poincar\'e inequality.
% The variance--surface-area inequality instead relates boundary size to volume.
These inequalities have a wide range of applications \cite{FriedgutKalai96,GG11,KoehlerLifshitzMinzerMossel23,khot2018monotonicity}, such as sharp-threshold phenomena and voting models with correlated preferences.
% For Boolean functions encoding graph properties, such as the existence of a Hamiltonian cycle, these inequalities characterize the robustness of the property to local edge changes.
